
\section{Introduction}
\label{sec:1}

\subsection{Background and Motivation}

The development of wireless communication
technologies has created significant opportunities
for improving transportation safety, efficiency,
and connectivity. Intelligent Transportation
Systems (ITS) integrate communication, sensing,
and information processing technologies to
support safer and more efficient road networks.

Vehicular Ad Hoc Networks (VANETs) represent
an important component of ITS. A VANET is
a specialized form of Mobile Ad Hoc Network
(MANET) in which vehicles communicate with
other vehicles and roadside infrastructure.
These networks support the exchange of
traffic-related information without requiring
all communication to pass through centralized
infrastructure \cite{Hartenstein2008}.

The primary motivation behind VANET research
is the need for timely and reliable
information exchange in highly dynamic
transportation environments. Vehicles may
share information about emergency braking,
traffic congestion, road hazards, and
unexpected driving conditions.

Such communication can improve situational
awareness and support cooperative
transportation applications. However,
the practical effectiveness of these
services depends on wireless connectivity,
communication reliability, and the
accuracy of exchanged information
\cite{Zeadally2012}.

Unlike conventional wireless networks,
VANETs operate under conditions
characterized by high mobility,
rapid topology changes, and
variable vehicle density.
These characteristics introduce
significant challenges for
network design and communication
protocols \cite{Karagiannis2011}.

\subsection{VANET Communication Environment}

Vehicular communication commonly includes
Vehicle-to-Vehicle (V2V) and
Vehicle-to-Infrastructure (V2I) communication.

V2V communication enables nearby vehicles
to exchange information directly.
For example, a vehicle detecting a
hazardous road condition may distribute
a warning to surrounding vehicles.

V2I communication allows vehicles
to interact with Roadside Units (RSUs)
and other transportation infrastructure.
These connections may support
traffic information distribution,
roadside services, and
infrastructure-assisted applications.

Both communication modes form part
of the broader Vehicle-to-Everything
(V2X) concept, which also includes
communication with pedestrians
and external communication networks.

Vehicular networking architectures
must consider the requirements
of different applications.
Safety-related services may require
timely message dissemination,
while traffic management and
information services may prioritize
coverage, connectivity, or
data capacity \cite{Karagiannis2011}.

\subsection{Research Problem}

Despite their potential benefits,
VANETs face technical challenges
that distinguish them from
traditional wireless networks.

One major challenge is
rapidly changing network topology.
As vehicles move, communication
links may be established or
disconnected within short
periods of time.

Consequently, routing mechanisms
must adapt to changing
connectivity while limiting
communication overhead
and transmission delay.

Another challenge is
the variation in vehicle density.
Sparse traffic conditions
may create disconnected
network segments, while
dense traffic can increase
wireless channel contention
and packet collisions.

Wireless propagation conditions,
road geometry, and physical
obstacles can also influence
communication reliability.

Furthermore, VANETs must
address security and privacy
concerns. Malicious participants
may attempt to inject
false information,
disrupt communication,
or exploit vehicle
location information.

These challenges demonstrate
that reliable vehicular
communication requires
careful consideration
of mobility, routing,
wireless access,
security, and
application requirements
\cite{Zeadally2012,Karagiannis2011}.

\subsection{Research Objectives}

The main objective of this paper
is to provide a structured
technical review of
Vehicular Ad Hoc Networks
from a wireless networking
perspective.

The specific objectives are:

\begin{enumerate}
\item To examine the fundamental
architecture of VANETs and
the functions of On-Board Units
(OBUs) and Roadside Units (RSUs).

\item To explain V2V, V2I,
and broader V2X communication
mechanisms and their applications.

\item To review major routing
protocol categories and discuss
their advantages and limitations
under vehicular mobility.

\item To investigate important
road safety, traffic management,
and information services.

\item To identify challenges
related to connectivity,
communication reliability,
security, and privacy.

\item To review established
simulation frameworks and
performance evaluation metrics.

\item To discuss important
technical trade-offs and
future research directions.
\end{enumerate}

\subsection{Scope and Review Methodology}

This paper adopts a
literature-based technical
review approach.

Its scope includes
VANET architectures,
vehicular communication
technologies, routing
protocols, applications,
security considerations,
and performance evaluation.

The discussion draws on
published survey research,
technical studies, and
established open-source
simulation frameworks.

Previous surveys provide
a foundation for understanding
VANET requirements,
architectures, challenges,
and communication approaches
\cite{Hartenstein2008,Karagiannis2011,Zeadally2012}.

Simulation is an important
research methodology
because it enables
communication behavior
to be studied under
controlled mobility
and network conditions.

For example, the Veins
framework integrates
road traffic simulation
with wireless network
simulation, enabling
coordinated investigation
of vehicle movement
and communication
behavior \cite{Sommer2011}.

This paper reviews
existing simulation
approaches and
performance metrics.
It does not claim
original simulation
experiments or
newly measured
performance results.

\subsection{Contributions of This Paper}

The contribution of
this paper is a
structured synthesis
of important VANET
research topics.

First, it examines
the relationship
between network
architecture,
communication
technologies,
and transportation
applications.

Second, it reviews
routing approaches
and discusses their
limitations under
highly dynamic
mobility conditions.

Third, it considers
security, privacy,
and communication
reliability as
interconnected
design challenges.

Finally, it reviews
simulation frameworks
and performance metrics
and discusses how
experimental assumptions
can influence the
interpretation of
VANET research.

\subsection{Organization of the Paper}

The remainder of this paper is organized
as follows.

Section~II introduces VANET architecture,
system components, and vehicular
communication technologies.

Section~III discusses routing protocols,
vehicular applications, and major
technical challenges.

Section~IV reviews simulation frameworks
and performance evaluation metrics.

Section~V presents a comparative analysis
of published studies and routing approaches.

Section~VI examines modern V2X communication
technologies, including IEEE 802.11p,
LTE-V2X, IEEE 802.11bd, and 5G NR-V2X.

Section~VII discusses open research
challenges and future directions,
including security, privacy,
scalability, and emerging approaches.

Finally, Section~VIII presents the
conclusions and future work.
